\begin{center}\large\bfseries Abstract\end{center}
\noindent\textit{Distributionally robust optimization (DRO) addresses uncertainty about a return law in addition to risk under that law. This paper reviews established DRO methods and evaluates a reproducible portfolio pipeline without claiming a new optimization method. A persistence-anchored residual LSTM predicts changes in log rolling volatility. Physical Heston parameters are fitted annually to past returns using a latent-variance particle likelihood. A four-ETF mean-CVaR problem combines a Wasserstein penalty with quadratic risk. A preliminary single-asset analysis over 2019--2024 motivated the design, so that period is exploratory; 2025 supplies a separate retrospective confirmation period. Residual forecasting lowers RMSE against persistence by 16--21\% in exploration and 22--29\% in confirmation, although a simple rolling-window update remains competitive. Fitted Heston models have mixed predictive likelihood gains. A matched 20,000-scenario comparison under three independent seeds reduces Monte Carlo variation below the average robustness tilt, while robust-versus-nominal return differences remain statistically inconclusive. Strong 2025 performance is interpreted in light of a gold rally and concentrated gold holdings in a historical benchmark. A cash-enabled comparison separates asset composition from de-risking. The public release provides code, model records, aggregate results, and input hashes; numerical replication is conditional on access to matching provider inputs.}
\par\vspace{7pt}
\noindent\textbf{Keywords:} distributionally robust optimization; stochastic volatility; residual learning; Heston calibration; conditional value-at-risk; multi-asset portfolios.

\section{Introduction}
Portfolio decisions depend on an estimated distribution of future returns. Risk criteria summarize losses under that distribution, whereas distributionally robust optimization (DRO) evaluates specified perturbations of the distribution itself. A flexible volatility model and a conservative decision criterion address different parts of this problem. Combining them is therefore a modeling choice to validate, rather than evidence of improved investment performance.

\begin{figure}[t]
\centering
\begin{tikzpicture}[node distance=4mm,box/.style={draw,rounded corners=1pt,align=center,text width=6.55cm,inner sep=4pt,font=\small},half/.style={draw,align=center,text width=3.0cm,inner sep=3pt,font=\small},arr/.style={-{Latex[length=2mm]},thick}]
\node[box] (a) {Information after close $t$\\Past ETF prices and VIX};
\node[half,anchor=north] (b) at ([xshift=-17mm,yshift=-5mm]a.south) {Residual forecast\\Persistence plus\\learned change};
\node[half,anchor=north] (c) at ([xshift=17mm,yshift=-5mm]a.south) {Physical Heston model\\Past-return particle\\likelihood};
\node[box,below=24mm of a] (d) {Correlated multi-asset scenarios\\Common estimated drifts and price-shock dependence};
\node[box,below=of d] (e) {Convex mean--CVaR allocation\\Wasserstein penalty and quadratic risk};
\node[box,below=of e] (f) {Next-open four-ETF holdings\\Daily wealth, trading costs, and diversification};
\draw[arr] (a)--(b);\draw[arr] (a)--(c);\draw[arr] (b)--(d);\draw[arr] (c)--(d);\draw[arr] (d)--(e);\draw[arr] (e)--(f);
\end{tikzpicture}
\caption{The computational pipeline. Forecasting, parameter estimation, and optimization have separate benchmarks and diagnostics. The primary portfolio is fully invested, which holds gross exposure constant across model variants.}
\label{fig:pipeline}
\end{figure}

Mean--variance analysis organizes portfolio choices through expected returns and covariances \cite{markowitz}. Its formulation does not require normal returns, although those moments need not capture all downside preferences. CVaR makes tail losses explicit \cite{rockafellar}, while a Wasserstein ambiguity set specifies a geometry for changes in outcomes \cite{esfahani}. Stochastic-volatility dynamics describe time-varying dispersion; they do not eliminate distributional misspecification.

This paper develops a study project on these established components as a selective review and reproducible experiment. Three design choices are central. First, a neural forecast should demonstrate value beyond the persistence already present in its target. Second, stochastic-volatility parameters should be estimated from historically available data rather than treated as universal constants. Third, a portfolio experiment should distinguish diversification from a scalar switch between equity and cash.

We address these questions with a residual forecast, a fitted latent-variance model, and a four-asset convex allocation problem. The forecast includes an identity path, a zero-initialized correction head, and a validation-only shrinkage rule. The Heston likelihood conditions variance transitions on observed returns, retaining leverage. The allocation includes an explicit positive-definite quadratic risk term and uses a Euclidean transport metric, so robustness depends on the vector of weights. No per-asset cap forces diversification.

A preliminary single-asset analysis over 2019--2024 motivated the design, so that period is labeled exploratory. Economic coefficients were fixed before the multi-asset outcomes were inspected. A separately specified 2025 sample is a retrospective confirmation period, not a publicly preregistered prospective trial. Numerical accuracy is assessed through tighter likelihood searches and matched 20,000-scenario allocation runs under three independent seeds. Those checks do not select economic settings by future performance.

The main findings are qualified. Residual prediction substantially improves the neural comparison with persistence, but a mechanical rolling-window baseline explains much of the available target predictability. Estimated Heston parameters vary across assets and years, yet their predictive likelihood gains are mixed. Multi-asset allocations permit diversification at fixed exposure. A higher-scenario comparison reduces simulation variation, while incremental DRO return effects remain small and statistically uncertain. These findings support an educational synthesis rather than an algorithmic novelty or universal performance claim.
